\chapter{Demostraciones de paradigmas de control}
Este capítulo presenta la colección de aplicaciones Flutter de demostración desarrolladas como contribución de referencia del proyecto al DEE. Cada demostración explora un paradigma específico de control o monitorización del dron de forma autónoma. El objetivo es proporcionar un ejemplo mínimo, bien documentado y reutilizable para cada tecnología, que pueda ser estudiado, adaptado e integrado de forma independiente por futuros desarrolladores del ecosistema.

El apartado 4.1 describe la organización del repositorio. Los apartados 4.2 a 4.X presentan cada demostración individualmente: joysticks virtuales, control por voz, control inercial mediante IMU, localización GPS en mapa y streaming de vídeo WebRTC con detección de objetos.

\section{Organización y estructura del repositorio}
Todas las demostraciones se organizan como proyectos Flutter independientes, dentro de la carpeta “demostraciones/” del repositorio del proyecto. Cada proyecto dispone de su propio “pubspec.yaml” y de su propio “main.dart” como punto de partida. Esta independencia es un requisito de diseño elegido, ya que permite que un desarrollador que solo necesite entender el control por voz, por ejemplo, pueda clonar el proyecto, abrir la carpeta que toca y ejecutar la demostración sin necesidad de configurar el resto del sistema.

Todos los ficheros de código están comentados, documentando no solo el funcionamiento sino también como debería adaptarse para conectarse a un dron real. Este enfoque orientado a la reutilización es lo que distingue esta colección de un conjunto de prototipos, y la convierte en un recurso de referencia para la comunidad.

\section{Catálogo de widgets}

\subsection{Caso de uso}
Una estación de control usa una paleta reducida, pero muy específica de componentes de interfaz. Botones que cambian de color según el estado, controls que responden a un pulsado contínuo, sliders de zoom, selectores de modo, listas reordenables etc. Cada uno de estos widgets tiene un comportamiento, unos estados visuales y unas características que no són evidentes a primera vista y que, si se implementan incorrectamente pueden causar varios problemas. 

Esta demostración es un poco distinta a las que le proceden. No explora un paradigma de control sino que actúa como catálogo interactivo y ejecutable de los 10 tipos de widget principales empleados en EZDrone, con cada igual a como se encuentran en este, con sus variantes visuales documentadas y una nota que explica los detalles de implementación relevantes. El objetivo es que un desarrollador que quiera reutilizar algún componente pueda ver el widget en acción, entender su lógica y copiarlo.

\subsection{Implementación}
La demostración se estructura como una pantalla de scroll vertical con diez secciones, cada una encabezada por el nombre del widget, su uso en EZDrone y una nota que documenta los aspectos no evidentes de su implementación.

El primer widget es ElevatedButton.icon, el botón de acción principal usado en todas las pantallas de EZDrone. La sección replica la barra de controles de vuelo (arm, take off, land, rtl y disconnect). Cada botón tiene habilitación condicional: arm solo está disponible cuando isConnected es verdadero y el dron no está ya armado ni volando, y así. Al pulsar los botones en la demo se sigue el ciclo real de un dron, cosa que permite ver visualmente su correcto funcionamiento.

El segundo widget es GestureDetector combinado con AnimatedContainer, patrón que usar los botones del HUD de la pantalla Tello, en los botones de flip y en los chips de modo. La sección muestra tres botones con estados enabled/disableed y active/inactive. AnimatedContainer solamente anima el color de fondo y el grosor del borde, haciendo que transicione suavemente entre estados.

El tercer widget es el botón de icono compacto (GestureDetector sobre un Container), usado para los botones de captura de foto, inicio/parada de grabación y cambio de cámara y mapa. A diferencia del widget anterior no usa AnimatedContainer, ya que el cambio de estado en estos botones es instantáneo. La sección muestra el botón de grabación con un toggle, donde al pulsarlo, cambia el icono y el color.

El cuarto widget documenta el patrón de presión continua, mediante GestureDetector con onLongPressStart, onLongPressEnd, onTapDown, onTap y onTapCancle, empleado en los botones de ajuste de la altitud en modo IMU. La combinación de los cinco callback cubre todos los posibles casos. 

El quinto widget es el Slider de zoom de la cámara, con el tema personalizado que usa SliderThemeData. La sección muestra el slider funcional con la etiqueta actualizada según el zoom.

El sexto widget es el DropdownButton del selector de modo de detección YOLO. La nota documenta los parámetros necesarios para que el desplegable encaje visualmente en la superposición sobre el vídeo.

El séptimo widget es el chip de modo, ModeChip y ConnectionChip, con AnimatedContainer. La sección muestra los selectores de modo de control (Classic, Voice, IMU) y de modo de dron (ArduPilot, SITL, Tello) completamente funcionales, con transición entre estados.

El octavo widget es el ModalBottomSheet combinado con DraggableScrollableSheet, usado en la pantalla de ayuda y en el selector de modos de la pantalla Tello. La sección incluye un botón que abre el sheet, mostrando el comportamiento de arrastre con scroll.

El noveno widget es el AlertDialog de confirmación. En este caso para la confirmación de eliminación.

El décimo widget es el ReordenableListView con handles de arrastre manuales mediante ReordenableDragStartListener, usado en la FlightPlanScreen en la lista de waypoints. La sección muestra una lista de cuatro waypoints de ejemplo completamente reordenables, con borrado individual. La nota documenta el detalle crítico del reordenamiento, que es necesaria por la forma en que Flutter calcula los índices tras eliminar el elemento de su posición original.

\subsection{Integración en EZDrone}
Esta demostración no se integra en EZDrone como funcionalidad, sino que actúa como agrupación de los diferentes elementos de diseño de la aplicación. Todos los widgets documentados se usan directamente en EZDrone con los mismos parámetros, colores y comportamientos descritos.

\subsection{Demostración}

\section{Joystick viruales}

\subsection{Caso de uso}

El mando RC de doble palanca, es el sistema de control de dron más usado entre los pilotos. Control cuatro canales de forma independiente: throttle y yaw en la palanca izquierda, pitch y roll en la derecha. Replicar esta disposición sobre una pantalla táctil convierte un móvil en un mando RC sin ningún tipo de hardware adicional.

El objetivo de esta demostración es documentar cómo se comportan los joysticks virtuales, rango de salida, comportamiento al soltarlo o mapeo a los cuatro ejes de vuelo, antes de integrarlo en una aplicación completa. Entender cómo funcionan esos valores es imprescindible para poder hacer una correcta conversión a señales PWM. 

\subsection{Implementación}
La demostración utiliza el paquete flutter-joystick (versión 0.0.4), que proporciona un widget “Joystick” que rastrea la posición del dedo relativamente al centro, y expone los valores continuos X e Y mediante un callback “listener”. La interfaz presenta dos joysticks configurados con “mode: JoystickMode.all”, que permite movimiento libre en los 360º, siguiendo la disposición de un control RC estándar.

Cada eje produce un valor en el rango [-1.0, 1.0], donde 0.0 representa la posición de reposo central. Una característica esencial del widget, es que cuando sueltas el dedo, el stick vuelve automáticamente al centro y el listener emite (0.0, 0.0). Esta característica es fundamental para el contexto de control de drones, sobre todo en modos como son el ALT-HOLD o el LOITER, donde un valor cero significa detener el movimiento.

Los valores de los cuatro ejes se muestran en pantalla en tiempo real, para que el desarrollador pueda verificar el rango de salida y la dirección antes de implementarlos en una aplicación. La conversión de estos valores normalizados a una señal PWM se realiza mediante la siguiente fórmula:

PWM = 1500 + [valor] * 500

Esta expresión mapea un valor de -1.0 a 1000 µs (deflexión máxima izquierda o abajo), 0.0 a 1500 µs (posición neutra) y +1.0 a 2000 µs (deflexión máxima derecha o arriba), cubriendo exactamente el rango PWM estándar de señales RC que espera el autopiloto.

\subsection{Integración en EZDrone}
En EZDrone, el mismo widget Joystick y la misma fórmula de conversión se utilizan en la pantalla de vuelo clásica (FlightScreen). La diferencia respecto a la demostración es la capa de transmisión. En lugar de mostrar los valores en la pantalla, los cuatro ejes se codifican en formato JSON ( {lx, ly, rx, ry} ) y se envían a la estación de tierra Python a través de un DataChannel WebRTC cada 50 ms (20Hz). La estación tierra recibe este paquete en el callback “on-joystick”, convierte los valores a PWM usando la misma fórmula y los envía al autopiloto mediante “dron.send-rc(roll, pitch, throttle, yaw)”.

Sobre la demostración base, EZDrone añade dos extras, el primero es una zona muerta de  ±0.1, donde los valores del joystick que estén en esa zona, se mapean automáticamente a 1500 µs, evitando que pequeñas imprecisiones generen un movimiento indeseado en el dron. El segundo es una curva de exposición aplicada al eje de yaw, añadida para suavizar la respuesta en pequeños movimientos, ya que el giro en el eje yaw es muy sensible a pequeñas variaciones.

\subsection{Demostración}

\section{Control por voz}
\subsection{Caso de uso}
El control por voz representa un paradigma de control de drones que libera las manos del operador durante el vuelo. En escenarios de inspección o monitorización, el operador puede necesitar tomar notas, manipular equipos o simplemente mantener la vista fija en el dron sin apartar la atención de la pantalla. Un sistema de voz bien diseñado permite emitir comandos verbalmente con el móvil a un lado.

El reto principal es el diseño de una interfaz de voz para el control de drones es su robustez, un sistema que reaccione a cualquier palabra detectada, generará comandos falsos de forma constante en cualquier ambiente ruidoso. La solución adoptada en esta demostración es una arquitectura de wake word, donde el sistema escucha pasivamente todo el rato, pero solo procesa comandos tras detectar la palabra de activación “dron”. Esto establece un umbral de activación sin requerir que el operador pulse ningún botón.

\subsection{Implementación}
La demostración utiliza el paquete speech-to-text (versión 6.6.2), que actúa como puente hacia el motor de reconocimiento de voz nativo del dispositivo: Google Speech en Android y el motor en Apple en iOS y macOS. Todo el procesamiento de reconocimiento es local, sin llamadas APIs externas, cosa que hace que se elimine la latencia de la red, y permite el funcionamiento sin conexión a internet una vez cargado el modelo.

El sistema de reconocimiento se modela como una máquina de estados de cuatro estados, implementada mediante el enum ListenState:

\begin{itemize}
    \item \textbf{idle} El sistema está inactivo, y el micrófono  apagado.
    \item \textbf{waitingWakeWord}  El micrófono está activo, esperando a que se diga la palabra de activación “dron”. 
    \item \textbf{waitingCommand}  La wake word ha sido detectada y se está esperando algún comando de vuelo.
    \item \textbf{commandLocked} El comando se ha ejecutado con éxito y el sistema permanece activo indefinidamente para recibir nuevos comandos.
\end{itemize}

Esta máquina de estados controla la lógica de reconocimiento y el estado visual de la interfaz, de forma que el operador siempre tiene una referencia visual de en qué estado se encuentra el sistema en todo momento.

El bucle de escucha se mantiene activo de forma continua mientras el sistema no está en estado idle. El motor de reconocimiento se reinicia automáticamente en los callback onStatus y onError para evitar que un silencio o un error puntual interrumpa la escucha. El parámetro pauseFor se establece en 4 segundos, para evitar que el motor corte comandos compuestos como podrían ser “mover adelante” o “volver al despegue”. El idioma de reconocimiento se puede cambiar, pero en esta demostración se fija en español “es-ES”.

Para gestionar el caso en que el operador olvida desactivar el sistema, se implementa un temporizador de apagado automático de 10 segundos. Este temporizador se resetea cada vez que se detecta cualquier voz y se cancela permanentemente cuando se ejecuta un comando con éxito, haciendo que el sistema se quede activo mientras el operador lo está usando.

Los comandos reconocidos se implementan en el método -identifyCommand, y cubren el ciclo de vuelo completo:

"armar"             -	ARM

"despegar"       	-	TAKEOFF

"subir"		        -	Ajuste de altitud positivo

"bajar" 	        -	Ajuste de altitud negativo

"mover adelante" 	-	Navegación hacia adelante

"mover atrás"  	    -	Navegación hacia atrás

"mover derecha" 	-	Navegación lateral derecha

"mover izquierda"	-   Navegación lateral izquierda

"aterrizar"  	    -   LAND

"volver a despegue" -   RTL

\subsection{Integración en EZDrone}
La integración del control por voz en EZDrone introduce una diferencia arquitectónica respecto a la demostración, ya que Flutter Web no puede utilizar el paquete speech-to-text (porque este depende de APIs nativas de la plataforma que no están disponibles en navegador). La solución que se adopta es la Web Speech API, soportada por todos los navegadores, que expone un puente  de interoperabilidad de JavaScript a Dart.

Esta capa se implementa en la clase WebSpeech, ubicada en el frontend Flutter. La clase define tipos de extensiones Dart, sobre los objetos JS correspondientes, lo que permite manipularlos desde Dart. 

En EZDrone la máquina de estados de cuatro fases de la demostración se sustituye por un modelo push-to-talk, más directo y adecuado para el uso de control de un dron. El operador mantiene pulsado el botón de micrófono, dicta el comando y lo suelta. Al soltar se llama al comando -speech.stopListening( ), y el texto reconocido se procesa. Este modelo elimina la necesidad de usar una wake word, ya que solo se dictan comandos cuando el operador lo decide.
\subsection{Demostración}

\section{Control inercial mediante IMU}

\subsection{Caso de uso}
La mayoría de los teléfonos móviles tienen una unidad de medida inercial (IMU) compuesta por acelerómetro, giroscopio y magnetómetro, cuya salida describe la orientación física del dispositivo en un espacio tridimensional. Mapear esta orientación a los ejes de control del dron crea una forma de controlarlo intuitiva, ya que inclinar el teléfono hacia adelante hace que el dron avance, inclinarlo hacia lateralmente lo desplaza hacia un lado y así. El operador sujeta el dispositivo con una empuñadura natural, y el control se genera mediante la postura y no mediante movimientos precisos de dedo.
Este paradigma es especialmente interesante para operadores menos familiarizados con la disposición de joysticks RC.

El reto técnico principal es que Flutter Web no puede acceder a los sensores del dispositivo directamente. El único camino disponible para ello es usar la API DeviceOrientationEvent del navegador, que requiere contexto HTTPS y en iOS, una solicitud explícita de permiso vinculada a un gesto del usuario.

\subsection{Implementación}
La solución se centra en torno a un puente JavaScript definido en el fichero index.html. Este puente registra un listener sobre el evento deviceorientation que va almacenando continuamente los tres ángulos de Euler en variables globales, y expone cinco funciones al entorno: “getAlpha( )” yaw, rango [0º, 360º], “getBeta( )”: pitch, rango [-180º, 180º],  “getGamma( )” roll, rango [-90º, 90º],  requestMotionPermission( ), para la compatibilidad con iOS, y finalmente stopOrientation( ), que elimina el listener.

En el lado de Dart, estas funciones se declaran mediante la anotación @JS, que permite llamar a funciones externas:

@JS('getBeta')  external double getBeta();
@JS('getGamma') external double getGamma();
@JS('getAlpha') external double getAlpha();

Un timer de 50 ms, llama continuamente a estas funciones, y se actualizan las variables internas de pitch, roll y yaw. Estos valores se muestran en tiempo real en tres recuadros. El botón “Activar sensores” aprovecha para combinar el inicio del timer con la petición obligatoria que requiere iOS.

Un aspecto práctico documentado en la demostración es que DeviceOrientationEvent requiere HTTPS, mientras que el servidor de Flutter sirve sobre HTTP. Para probar la demostración en un dispositivo físico, es necesario exponer el servidor local a través de un túnel HTTPS, la demostración incluye instrucciones paso a paso para hacerlo mediante ngrok.

\subsection{Integración en EZDrone}
En EZDrone, el mismo puente JavaScript, y el mismo timer se utilizan directamente en ImuFlightScreen. La integración añade dos elementos respecto a la demostración base.

El primero es el tratamiento de zona muerta.  Los ángulos brutos de los sensores no se pueden enviar directamente, ya que incluso con el dispositivo quieto, los valores fluctúan ligeramente, lo que generaría movimientos continuos. La función normalForward( double beta) implementa una zona muerta (entre -5º y 15º) para el eje del pitch, dado que el dispositivo en posición de reposo natural suele estar en esa inclinación. 
La función normalLateral (double gamma), hace lo mismo para el movimiento lateral con una zona muerta de ±10º. Una vez fuera de la zona muerta, se aplica una rampa lineal de 40º donde se mapean los valores para poder transformarlos a velocidades.

El segundo es el bloqueo de orientación de la pantalla. Cuando el dispositivo se inclina para controlar el dron, el sistema operativo puede girar la interfaz al detectar que el dispositivo se ha girado de más, haciendo la pantalla de EZDrone inutilizable. Esto se soluciona llamando a sceen.orientation.lock( ) mediante el puente JS definido en html.index. Al salir de la pantalla, unlockOrientationEZ( ) libera este bloqueo.

\subsection{Demostración}

\section{Localización GPS en mapa}

\subsection{Caso de uso}
Conocer la posición del operador en tiempo real es una funcionalidad con mucho valor operativo, permite ver la distancia real entre el operador y el dron, mostrar ambas posiciones en el mismo mapa, y en un escenario futuro, implementar modos de seguimiento o establecer zonas de exclusión relativas a la posición del operador. 

Esta demostración documenta cómo obtener la posición GPS del dispositivo y visualizarla en tiempo real sobre un mapa de código abierto.

\subsection{Implementación}
La demostración utiliza dos paquetes: geolocator (versión 13.0.2) para la obtención de la posición GPS, y flutter-map (versión 7.0.2) para el renderizado del mapa, con latlong2 (versión 0.9.1) para el manejo de coordenadas. La elección de flutter-map sobre alternativas es intencional, ya que no requiere claves API y es completamente abierta. Utiliza tiles de OpenStreetMap, lo que la hace adecuada para el contexto del proyecto, y elimina dependencias de servicios de pago.

La posición se obtiene mediante Geolocator.getPositionStream( ), que abre un stream que emite actualizaciones cada vez que el dispositivo se desplaza al menos 1 metro (configurable). Cada posición nueva se acepta siempre si es la primera (para evitar un bucle de espera) o si la precisión es inferior a 50 metros. Las lecturas con mayor precisión se descartan para evitar saltos bruscos en el mapa. 
La precisión GPS actual se muestra en un badge superpuesto sobre el mapa con una codificación de colores, verde por debajo de 5 m, naranja por debajo de 20 m y rojo por encima.

\subsection{Integración en EZDrone}
En EZDrone la misma lógica de stream GPS se usa en los métodos startUserLocation( ) y stopUserLocation( ) del DronProvider. La posición del operador se almacena en userPosition, y se inicia automáticamente al conectar al dron, deteniéndose al desconectar. La posición se visualiza dentro del mapa como un icono de teléfono móvil azul, permitiendo al operador percibir de forma visual la separación entre él y el dron.

\subsection{Demostración}

\section{Streaming de vídeo WebRTC con detección YOLO}

\subsection{Caso de uso}
La visión en primera persona (FPV) es una capacidad fundamental en cualquier estación de control de drones para misiones de inspección, fotografía aérea o navegación en entornos complejos. Esta demostración documenta cómo establecer un stream de vídeo en tiempo real desde una cámara conectada al ordenador. Se utiliza WebRTC como protocolo de transporte, y se integra detección de objetos en tiempo real mediante YOLO sobre el stream antes de su transmisión.

\subsection{Implementación en el emisor Python}
El emisor es un script Python, que utiliza las bibliotecas aiortc y websockets. Define una clase CustomVideoStreamTrack que hereda de VideoStreamTracl de aiortc y sobreescribe el método asíncrono recv( ).

En cada llamada a recv( ) (que es en cada frame que se tiene que transmitir), el track captura un frame de la cámara mediante OpenCV, ejecuta interferencia YOLO cada 25 frames, dibuja los bounding boxes sobre el frame según el modo de detección activo, añade un timestamp y convierte el frame al formato rgb24 requerido por aiortc.

El servidor WebSocket permanece abierto durante toda la sesión, actuando tanto de servidor de señalización como canal de control en tiempo real para los comandos que Flutter envía durante el streaming (cambios de modo de detección: all/person/none, solicitud de captura de frame e inicio/parada de grabación).

\subsection{Implementación en el cliente Flutter}
El cliente Flutter utiliza la librería flutter-webrtc. Al conectar, crea una RTCPeerConnection con el servidor STUN de Google, registra un callback onTrack que asigna el stream de vídeo recibiendo al renderer RTCVideoRenderer, y abre el canal WebSocket hacia el emisor Python.

La captura de pantalla y la grabación de vídeo se realizan en el cliente mediante interoperabilidad JS, sin enviar ningún dato al servidor Python. Las funciones captureDroneFrame( ) y startDroneRecording( ) / stopDroneRecording( ), definidas en el index.html, acceden al stream y lo extraen mediante captureStream( ) y utilizan la API MediaRecorder del navegador para grabar en formato WebM. Dart llama a estas funciones mediante @JS y la biblioteca dart:js-interop.

\subsection{Integración en EZDrone}
En EZDrone, el stream WebRTC se recibe en WebRTCLogic y el MediaStream resultante se expone al DronProvider mediante el callback onRemoteStream. La pantalla de vuelo clásica puede alternar entre el mapa satelital y el stream de vídeo mediante el botón de intercambio, renderizando el stream mediante el widget RTCVideoView. Las mismas funciones JS de captura y grabación se invocan desde el DronProvider mediante jsCapture( ) y jsStartRecording( ) / jsStopRecording( ), exponiendo estas funciones como botones en la interfaz de vuelo. El modo de detección YOLO se controla desde la misma pantalla mediante un widget dropdown que deja elegir entre personas, todo y nada. Entonces, mediante el topic MQTT mobileFlutter/groundStation/detectionMode, Flutter envía la elección a la estación tierra, donde se produce el cambio de detection-mode.

\subsection{Demostración}

\section{Planificación de misiones autónomas}

\subsection{Caso de uso}
Los paradigmas de control descritos en los apartados anteriores son todos reactivos, es decir, el operador emite comandos en tiempo real y el dron responde a ellos de forma inmediata. Existen una clase de operaciones donde lo que se busca es definir la misión con antelación y dejar su ejecución al autopiloto, sin requerir de la atención e intervención continua del operador (como podrían ser inspecciones de infraestructuras, cartografías aéreas, vigilancia de perímetros, riego de cultivos entre muchos otros).

ArduPilot soporta este modo de vuelo autónomo, permite que se le pasen una lista de waypoints con coordenadas GPS y altitudes, los almacena en la memoria y los ejecuta en modo AUTO.

Esta demostración abarca el problema de diseñar una interfaz de planificación de misiones accesible desde cualquier navegador. A parte de demostrar cómo construir una lista de waypoints de forma visual sobre un mapa, asignar acciones a cada uno de los waypoints, guardar los planes entre sesiones, exportarlos a un formato compatible con herramientas como MissionPlanner y enviarlos al dron para su ejecución.

\subsection{Modelos de datos}
La demostración describe el modelo de datos de la misión, que en EZDrone se implementa en tres clases distintas del Provider.

FlightWaypoint representa un punto de paso individual. Sus campos son lat y lon (coordenadas en grados), altM (altitud sobre el suelo en metros) y action, un objeto WaypointActoin. Este objeto contiene un tipo de acción a ejecutar al llegar al waypoint junto a su duración si es el caso. Los tipos de acciones disponibles son: none (volar sobre el punto sin detenerse), hover (mantener posición durante los segundos definidos), takePhoto (disparar la cámara), recordVideo (grabar durante el tiempo definido), rtl (retorno al punto de lanzamiento) y land (aterrizaje en el waypoint).

FlightPlan agrupa una lista de FlightWaypoints con un identificador único, un nombre editable y una marca de tiempo de creación. Esta clase también contiene un método para calcular las distancias entre waypoints consecutivos, totalDistanceM, usando la fórmula de Haversine: 

const R = 6371000.0;
final s = sin(dLat/2)*sin(dLat/2) + cos(lat1)*cos(lat2)*sin(dLon/2)*sin(dLon/2);
d += R * 2 * atan2(sqrt(s), sqrt(1 - s));

La elección de Haversine es más que correcta para esta escala, para misiones de varios kilómetros, la curvatura terrestre introduce errores no despreciables si se usa geometría plana.

\subsection{Implementación}
La demostración presenta una interfaz de planificación de misiones completamente funcional sin necesidad de dron ni simulador. 

La pantalla de planificación presenta una interfaz separada en dos zonas. La zona superior es un mapa satelital de flutter-map. Un tap sobre cualquier punto del mapa llama al callback de onTap, que añade un nuevo FlightWaypoint en el punto donde se ha pulsado. Los waypoints se ven como marcadores numerados, el primero en verde (indicando el inicio) y el último en rojo (indicando el final). Una línea lila los une en orden, y un badge superpuesto muestra el número de waypoints y la distancia total calculada.

La zona inferior es una lista de Waypoints, implementada usando ReordenableListView, que permite reordenar los puntos mediante arrastrar y soltar. Cada elemento de la lista muestra las coordenadas con cinco decimales, la altitud asignada, el chip de acción si existe y un botón para borrarlo. Al pulsar sobre un elemento se abre un AlertDialog de edición, donde el operador puede modificar la altitud y seleccionar el tipo de acción y su duración.

Los planes se guardan en el almacenamiento local del navegador, usando web.window.localStorage. La serialización y deserialización se realiza mediante los métodos toJson( ) y fromJson( ). 

La demostración incluye exportación a formato QGC, el estándar de misiones en el ecosistema de ArduPilot. Se genera un fichero de texto plano con todos los waypoints y su comando MAVLink correspondiente para sus acciones. La descarga se hace en el propio navegador usando la API Blob.

\subsection{Integración en EZDrone}
En EZDrone, la demostración se extiende con dos acciones adicionales, el envío del plan al dron y su ejecución. El botón UPLOAD serializa el plan como JSON y lo publica en el topic MQTT /uploadMission, la estación tierra lo almacena y confirma su recepción. Al pulsar START, la estación ejecuta el plan mediante dron.goto( ). Durante la ejecución, el mapa muestra el progreso con el trail y el waypoint activo resaltado comparado con los otros.

\subsection{Demostración}

\section{Registro y análisis de vuelos}

\subsection{Caso de uso}
El registro de los vuelos realizados tiene valor tanto operativo como de revisión. Conocer la trayectoria recorrida, la evolución de la carga de la batería, la velocidad pico alcanzada o el tiempo transcurrido desde el despegue, permite al operador evaluar el rendimiento del vuelo, detectar comportamientos anómalos y documentar las operaciones. Esta demostración aborda el diseño de una pantalla de historial de vuelos que combina visualización de trayectorias, reproducción temporal de la telemetría y exportación de los datos.

El objetivo de esta demostración es en particular doble, por un lado, documentar el modelo de datos y la lógica de visualización de forma aislada sin depender de un dron real, y por otro lado, servir como prueba antes de implementarlo en EZDrone con datos reales.

\subsection{Datos de ejemplo}
Dado que la demostración se ejecuta sin ningún tipo de dron o SITL, los datos han sido generados manualmente en el fichero sample-flight-data.dart. Este fichero contiene tres sesiones de ejemplo con coordenadas reales sobre el campus. La primera se trata de un vuelo cuadrado a 20 m en modo clásico (un tipo de modo de control de EZDrone) con SITL, de 4 minutos y 32 segundos, completado correctamente, la segunda es un vuelo triangular a 30 m en modo IMU (modo de control de EZDrone), de 6 minutos y 10 segundos, también completado. Finalmente la tercera se trata de un vuelo parcial a 15m en modo voz (modo de control de EZDrone), de 1 minuto y 48 segundos, interrumpido por desconexión. Esta combinación de sesiones con distintos modos de control, flags y estados de finalización, permite verificar diferentes opciones posibles de la interfaz.

Aparte, se ha implementado la función GenLog( ), capaz de generar la telemetría de cada sesión de forma coherente con la trayectoria. La altitud sube linealmente durante el primer 15/100 del vuelo, se mantiene con pequeñas fluctuaciones aleatorias, y en el último 15/100 baja. La batería decae de forma lineal, la velocidad es baja en los tramos de despegue y aterrizaje y se estabiliza entre 3 y 5.5 m/s en crucero. El rumbo en cada punto se calcula mediante trigonometría sobre el vector entre el punto actual y el siguiente, reproduciendo la orientación real del dron en cada segmento.

\subsection{Implementación}
La pantalla principal de la demostración, FlightLogScreen, presenta la lista de sesione en un ListView. Cada sesión se muestra en una tarjeta que incluye la fecha y hora de inicio, un chip de color con el modo de control usado, un chip naranja SITL cuando el vuelo ha sido simulado, un indicador de estado completado o interrumpido y cinco métricas calculadas: duración, distancia total, altitud máxima, desnivel y velocidad máxima.

Al pulsar sobre las tarjetas, se abre un bottom sheet, implementado mediante un DraggableScrollableSheet. La zona superior contiene un mapa con la trayectoria completa dibujada. Tres marcadores que identifican el punto de despegue (verde), el de aterrizaje / interrupción (rojo / naranja) y la posición actual del dron en cada momento.

El slider temporal en el elemento central de la bottom sheet. Su rango va de 0 al número de puntos del trail, y al arrastrarlo se actualizan simultáneamente el segmento visible del trail, la posición del dron y el texto de telemetría. 

Debajo del slider, se presenta una cuadrícula de ocho tarjetas que aparte de incluir las métricas de la tarjeta de resumen, incluyen la batería mínima registrada y el número total de snapshots del log. El botón DOWNLOAD CSV construye el fichero en formato CSV, y lo descarga en el navegador, pero en esta demostración se muestra el contenido en un diálogo si la descarga no está disponible.

\subsection{Integración en EZDrone}
En EZDrone, se usan directamente los datos obtenidos en el vuelo real o simulado, mediante el Provider. Al aterrizar, se cierra la sesión y se inserta en el flightHistory. La lista completa se guarda en el almacenamiento local. Y seguidamente se muestran en la pantalla de FlightLogScreen.

\subsection{Demostración}


